\documentclass[11pt]{article}

\usepackage[T1]{fontenc}
\usepackage{mathpazo}
\usepackage{microtype}
\usepackage[a4paper,margin=1.08in]{geometry}
\usepackage{amsmath,amssymb,amsthm,mathtools}
\usepackage[hidelinks]{hyperref}
\hypersetup{
  pdftitle={Infinite Projections in Full Group C*-Algebras of S-Arithmetic Groups},
  pdfauthor={Tseng},
  pdfsubject={Kazhdan projections, subgroup compression, and S-arithmetic groups}
}
\hyphenation{Kazhdan Proposition arithmetic}

\newtheorem{theorem}{Theorem}[section]
\newtheorem{proposition}[theorem]{Proposition}
\newtheorem{corollary}[theorem]{Corollary}
\newtheorem{lemma}[theorem]{Lemma}
\newtheorem{remark}[theorem]{Remark}

\newcommand{\C}{\mathbb C}
\newcommand{\F}{\mathbb F}
\newcommand{\Q}{\mathbb Q}
\newcommand{\R}{\mathbb R}
\newcommand{\Z}{\mathbb Z}
\newcommand{\Hh}{\mathcal H}
\newcommand{\Td}{\mathcal T}
\newcommand{\OK}{\mathcal O_K}
\newcommand{\OS}{\mathcal O_{K,\Sigma}}
\newcommand{\SL}{\operatorname{SL}}
\newcommand{\Cmax}[1]{C^*_{\max}(#1)}
\newcommand{\Cr}[1]{C^*_r(#1)}
\newcommand{\Cfd}[1]{C^*_{\mathrm{fd}}(#1)}

\title{Infinite Projections in Full Group $C^*$-Algebras\\
of $S$-Arithmetic Groups}
\author{Tseng}
\date{August 14, 2026}

\begin{document}
\maketitle

\begin{abstract}
Let $K$ be a number field, let $\Sigma$ be a nonempty finite set of nonzero
prime ideals of $\OK$, and let $n\geq4$.  We prove that
\[
 \Cmax{\SL_n(\mathcal O_{K,\Sigma})}
\]
contains an infinite projection and a proper isometry.  Hence it is not
finite and is neither stably finite nor MF, although
$\SL_n(\mathcal O_{K,\Sigma})$ is a finitely presented, residually finite
property-$(T)$ group.  In particular this applies to
$\SL_n(\Z[S^{-1}])$ for every nonempty finite set $S$ of rational primes.
More generally, if a countable discrete group $G$ contains a property-$(T)$
subgroup $\Gamma$ and some $t\in G$ satisfies
$t\Gamma t^{-1}\subsetneq\Gamma$, then the Kazhdan projection $p_\Gamma$
is infinite in $\Cmax G$; explicitly,
$u_{t^{-1}}p_\Gamma+(1-p_\Gamma)$ is a proper isometry.  To the best of our
knowledge, these are the first examples of countable discrete groups with a
non-finite---and hence non-MF---full group $C^*$-algebra.  The preparation of
this manuscript was AI-assisted.
\end{abstract}

\section{Introduction}

All groups in this paper are countable and discrete.  The faithful canonical
trace makes the reduced group algebra $\Cr G$ stably finite for every group
$G$.  By contrast, it was not known whether every full group $C^*$-algebra of
a discrete group must be finite.

A projection in a $C^*$-algebra is \emph{infinite} if it is
Murray--von Neumann equivalent to a proper subprojection of itself.  A unital
$C^*$-algebra is finite precisely when every isometry is unitary, and it is
stably finite if every matrix algebra over it is finite.  Choi asked whether
$C^*(G)$ can fail to be quasi-directly finite for a discrete group $G$
\cite[Question~1]{ChoiQDF}.  In the unital setting this asks whether $C^*(G)$
can contain a proper isometry.  Milh\o j and R\o rdam later recorded the
equivalent question whether a full group $C^*$-algebra can contain an infinite
projection \cite[Section~6]{MilhojRordam}.  The qualification ``discrete'' is
essential: nondiscrete locally compact examples failing quasi-direct finiteness
were already known \cite[Section~3]{ChoiQDF}.

An MF $C^*$-algebra embeds into a norm ultraproduct of matrix algebras and is
stably finite.  In March 2026, Shulman noted that no group was known whose
full or reduced group $C^*$-algebra is non-MF, and that stable finiteness of
full group algebras was not known in general \cite[Introduction]{Shulman}.

Our main examples are the standard higher-rank $S$-arithmetic groups.

\begin{theorem}[$S$-arithmetic groups]\label{thm:sarithmetic}
Let $K$ be a number field, let $\Sigma$ be a nonempty finite set of nonzero
prime ideals of $\OK$, and let $\OS$ be the ring of elements of $K$ integral
outside $\Sigma$.  For every $n\geq4$, the group
\[
 G=\SL_n(\OS),
\]
viewed as a discrete group, is finitely presented, linear, residually finite,
and has property $(T)$, while $\Cmax G$ contains a proper isometry and an
infinite projection.  Consequently, $\Cmax G$ is not finite and is neither
stably finite nor MF.
\end{theorem}

For $K=\Q$, write $\Z[S^{-1}]=\Z[p^{-1}:p\in S]$.  The subgroup and
compressor behind Theorem~\ref{thm:sarithmetic} are especially transparent:
for $p\in S$, put
\[
 \Gamma=
 \left\{
 \begin{pmatrix}A&x\\0&1\end{pmatrix}:
 A\in\SL_{n-1}(\Z),\ x\in\Z^{n-1}
 \right\},
 \qquad
 t=\operatorname{diag}(p,\ldots,p,p^{-(n-1)}).
\]
Then $\Gamma\cong\Z^{n-1}\rtimes\SL_{n-1}(\Z)$ has property $(T)$ and
\[
 t\Gamma t^{-1}
 =p^n\Z^{n-1}\rtimes\SL_{n-1}(\Z)
 \subsetneq\Gamma.
\]

The arithmetic theorem is an instance of the following general mechanism.

\begin{theorem}[Kazhdan compression criterion]\label{thm:criterion}
Let $G$ contain a property-$(T)$ subgroup $\Gamma$.  If some $t\in G$
satisfies
\[
 t\Gamma t^{-1}\subsetneq\Gamma,
\]
then the Kazhdan projection $p_\Gamma$ is infinite in $\Cmax G$, and
$\Cmax G$ contains a proper isometry.  Consequently, $\Cmax G$ is not
finite and hence is neither stably finite nor MF.
\end{theorem}

The mechanism is elementary.  In every representation, the Kazhdan
projection $p=p_\Gamma$ projects onto the $\Gamma$-fixed vectors.  The
subgroup inclusion implies
\[
 p u_{t^{-1}}p=u_{t^{-1}}p.
\]
Thus $v=u_{t^{-1}}p$ is an isometry in the corner $p\Cmax Gp$, with
\[
 v^*v=p,
 \qquad
 vv^*=u_{t^{-1}}pu_t\leq p.
\]
The quasi-regular representation on $\ell^2(G/\Gamma)$ detects that the last
inequality is strict.  Hence $p$ is infinite, and $v+(1-p)$ is a proper
isometry with initial projection $1$.

Every residually finite group is MF as a group
\cite[Proposition~2.13]{CarrionDadarlatEckhardt}.  Theorem
\ref{thm:sarithmetic} therefore shows, even among finitely presented
residually finite Kazhdan groups, that the converse to
\[
 \Cmax G\text{ is MF}\quad\Longrightarrow\quad G\text{ is MF}
\]
fails.

Wassermann used Kazhdan projections in full group $C^*$-algebras for
extension-theoretic obstructions \cite{Wassermann}.  Bekka obtained
finite-dimensional and tracial rigidity results for full group algebras of
arithmetic groups \cite{Bekka,BekkaSuper}; these do not imply failure of
finiteness.
One-sided subgroup compression also produces isometries in Hecke
$C^*$-algebras and semigroup crossed products \cite{LacaLarsen}.  Those
isometries live in Hecke completions rather than in the full group $C^*$-algebra
of the ambient discrete group.  Here property $(T)$ supplies the universal
fixed-vector projection inside the maximal group $C^*$-algebra itself.

Section~\ref{sec:compression} proves the criterion and locates its defect
projection.  Section~\ref{sec:arithmetic} proves the arithmetic theorem.
Section~\ref{sec:completion} records the resulting separation of maximal,
finite-dimensional, and reduced completions.

\section{The compression argument}\label{sec:compression}

We first recall two standard facts in the form needed below.

\begin{lemma}[Maximal subgroup inclusion]\label{lem:inclusion}
If $H<G$, the canonical map $\C[H]\to\C[G]$ extends to an injective
$*$-homomorphism
\[
 \Cmax H\longrightarrow\Cmax G.
\]
\end{lemma}

\begin{proof}
Restriction of representations gives $\|x\|_{\Cmax G}\leq
\|x\|_{\Cmax H}$ for $x\in\C[H]$.  Conversely, if $\pi$ is a unitary
representation of $H$, then the restriction to $H$ of
$\operatorname{Ind}_H^G\pi$ contains $\pi$ on the fibre over $H\in G/H$.
Hence
\[
 \|\pi(x)\|\leq\|\operatorname{Ind}_H^G\pi(x)\|
 \leq\|x\|_{\Cmax G}.
\]
Taking the supremum over $\pi$ proves the claim.
\end{proof}

\begin{proposition}[The Kazhdan projection]\label{prop:kazhdan}
If $\Gamma$ has property $(T)$, there is a central projection
$p_\Gamma\in\Cmax\Gamma$ such that, in every unitary representation $\pi$
of $\Gamma$, $\pi(p_\Gamma)$ is the orthogonal projection onto
$\Hh^\Gamma$.

More precisely, let $S\subset\Gamma$ be a nonempty finite symmetric
Kazhdan set with Kazhdan constant $0<\kappa\leq2$, and put
\[
 M=\frac1{|S|}\sum_{s\in S}u_s,
 \qquad
 L=\frac{1+M}{2},
 \qquad
 r=1-\frac{\kappa^2}{4|S|}<1.
\]
Then $L^j\to p_\Gamma$ in maximal norm and
\[
 \|L^j-p_\Gamma\|\leq r^j.
\]
\end{proposition}

\begin{proof}
Since $S$ is symmetric, $M=M^*$ and $\|M\|\leq1$, so $L$ is a positive
contraction.  If $\xi\perp\Hh^\Gamma$, then
\[
 \|\xi\|^2-\langle\pi(M)\xi,\xi\rangle
 =\frac1{2|S|}\sum_{s\in S}\|\pi(s)\xi-\xi\|^2
 \geq\frac{\kappa^2}{2|S|}\|\xi\|^2.
\]
Thus $\pi(L)$ is the identity on $\Hh^\Gamma$ and has norm at most $r$
on its orthogonal complement, uniformly over all representations.  It follows
that $L^j$ converges in $\Cmax\Gamma$ with the stated estimate, and its
limit has the required universal property.  The invariant-vector projection
commutes with $\pi(\Gamma)$ in every representation, so the limit is central.
See also \cite{DrutuNowak}.
\end{proof}

\begin{proof}[Proof of Theorem~\ref{thm:criterion}]
By Lemma~\ref{lem:inclusion}, regard $p=p_\Gamma$ as an element of
$A=\Cmax G$.  Put
\[
 v=u_{t^{-1}}p,
 \qquad
 q=u_{t^{-1}}pu_t.
\]
Let $U$ be any unitary representation of $G$, and let $P=U(p)$ be the
projection onto $\Hh^\Gamma$.  For $\xi\in\Hh^\Gamma$ and
$\gamma\in\Gamma$,
\[
 U_\gamma U_{t^{-1}}\xi
 =U_{t^{-1}}U_{t\gamma t^{-1}}\xi
 =U_{t^{-1}}\xi.
\]
Hence $U_{t^{-1}}\Hh^\Gamma\subseteq\Hh^\Gamma$, so
\[
 P U_{t^{-1}}P=U_{t^{-1}}P.
\]
Since this holds in every representation,
\[
 p u_{t^{-1}}p=u_{t^{-1}}p=v.
\]
Therefore $v\in pAp$ and
\[
 v^*v=p,
 \qquad
 vv^*=q\leq p.
\]

To prove strictness, consider the quasi-regular representation
\[
 \rho:G\curvearrowright\ell^2(G/\Gamma).
\]
The ranges of $\rho(p)$ and $\rho(q)$ are the $\Gamma$-fixed and
$(t^{-1}\Gamma t)$-fixed subspaces, respectively.  Since
$\Gamma\subsetneq t^{-1}\Gamma t$, choose
$h\in t^{-1}\Gamma t\setminus\Gamma$.  The vector $\delta_\Gamma$ is
$\Gamma$-fixed, whereas
\[
 \rho(h)\delta_\Gamma=\delta_{h\Gamma}\neq\delta_\Gamma.
\]
Thus $\rho(q)\neq\rho(p)$, and consequently $q<p$ in $A$.  This proves
that $p$ is infinite.  Finally, with
\[
 d=p-q,
 \qquad
 w=v+(1-p),
\]
the cross terms vanish and
\[
 w^*w=1,
 \qquad
 ww^*=1-d<1.
\]
Hence $w$ is a proper isometry.  An MF $C^*$-algebra is stably finite
\cite{BlackadarKirchberg}, so $A$ is not MF.
\end{proof}

The proof uses property $(T)$ only to obtain a projection whose image in
every representation of $G$ is the projection onto the $\Gamma$-fixed
vectors.  Whenever such a projection exists, the same argument applies.

The identities above also determine exactly where the defect $d=p-q$ is
visible.

\begin{proposition}[The compression defect]\label{prop:defect}
Under the hypotheses of Theorem~\ref{thm:criterion}, let $p,v,q,d$ be as
in its proof.
\begin{enumerate}
\item Every $*$-homomorphism from $\Cmax G$ to a finite unital
$C^*$-algebra annihilates $d$.  Hence every finite-dimensional
representation and every tracial GNS representation annihilates the closed
ideal generated by $d$.  In particular, $\Cmax G$ is not residually
finite-dimensional and has no faithful tracial state.
\item The elements $p,v,q,d$ lie in the kernel $I$ of
$\Cmax G\to\Cr G$.  The projection $p$ is infinite but not properly
infinite, and the $C^*$-algebra generated by $v$ in the corner with unit
$p$ is isomorphic to the Toeplitz algebra $\Td$, with defect projection
$d=p-vv^*$.
\item The quasi-regular representation $\lambda_{G/\Gamma}$ does not
annihilate $d$.  Consequently, the completion associated with
$\lambda_G\oplus\lambda_{G/\Gamma}$ is strictly larger than the reduced
completion, although it may coincide with the maximal completion.
\end{enumerate}
\end{proposition}

\begin{proof}
Let $\Phi:\Cmax G\to B$, with $B$ finite and unital.  Passing to the
finite corner $\Phi(1)B\Phi(1)$, assume that $\Phi$ is unital.  The partial
isometry $\Phi(v)$ implements
\[
 \Phi(p)\sim\Phi(q)\leq\Phi(p).
\]
Finiteness forces $\Phi(p)=\Phi(q)$, so $\Phi(d)=0$.  This covers
finite-dimensional representations.  For a tracial state $\tau$ one also
has $\tau(d)=\tau(v^*v)-\tau(vv^*)=0$; since $d$ is a projection, its
image in the GNS representation is zero.  As $d\neq0$, the asserted
failures of residual finite dimensionality and faithful traciality follow.

Strict compression forces $\Gamma$ to be infinite.  The restriction of the
left regular representation of $G$ to $\Gamma$ has no invariant vector, and
hence annihilates $p$.  Thus $p,v,q,d\in I$.  The relation $v^*v=p$ and
the strict inequality $vv^*<p$ make $p$ infinite.  The augmentation
character sends $p$ to $1$, so $p$ is not properly infinite.  Coburn's
theorem \cite{Coburn} identifies the algebra generated by the nonunitary
isometry $v$ in $p\Cmax Gp$ with $\Td$.

Finally, the proof of Theorem~\ref{thm:criterion} shows that
$\lambda_{G/\Gamma}(d)\neq0$, proving the last assertion.
\end{proof}

\section{\texorpdfstring{$S$}{S}-arithmetic groups}\label{sec:arithmetic}

\begin{proof}[Proof of Theorem~\ref{thm:sarithmetic}]
Set $m=n-1$.  Choose $\mathfrak p\in\Sigma$.  Since the ideal class group
of $K$ is finite, there are $h\geq1$ and $\alpha\in\OK$ such that
\[
 \mathfrak p^h=(\alpha).
\]
Thus $\alpha$ is a nonunit of $\OK$ and, since $(\alpha)=\mathfrak p^h$
is supported in $\Sigma$, a unit of $\OS$.

Inside $G=\SL_n(\OS)$ consider
\[
 \Gamma=
 \left\{
 \begin{pmatrix}A&x\\0&1\end{pmatrix}:
 A\in\SL_m(\OK),\ x\in\OK^m
 \right\}
 \cong \OK^m\rtimes\SL_m(\OK).
\]
By Borel--Harish-Chandra, applied after restriction of scalars from $K$ to
$\Q$, the diagonal archimedean embedding realizes $\Gamma$ as a lattice in
\[
 \prod_{v\mid\infty}
 \bigl(K_v^m\rtimes\SL_m(K_v)\bigr).
\]
See \cite{BorelHarishChandra}.
Each factor has property $(T)$ for $m\geq3$
\cite[Corollary~1.4.16]{BHV}, and property $(T)$ passes to lattices
\cite[Theorem~1.7.1]{BHV}.  Hence $\Gamma$ has property $(T)$.

Put
\[
 t=\operatorname{diag}(\alpha I_m,\alpha^{-m})\in G.
\]
For $A\in\SL_m(\OK)$ and $x\in\OK^m$, direct multiplication gives
\[
 t\begin{pmatrix}A&x\\0&1\end{pmatrix}t^{-1}
 =\begin{pmatrix}A&\alpha^{m+1}x\\0&1\end{pmatrix}
 =\begin{pmatrix}A&\alpha^n x\\0&1\end{pmatrix}.
\]
Consequently,
\[
 t\Gamma t^{-1}
 =\alpha^n\OK^m\rtimes\SL_m(\OK)
 \subsetneq\Gamma,
\]
where strictness follows because $\alpha$ is a nonunit.  Theorem
\ref{thm:criterion} gives the proper isometry and infinite projection.

For completeness, the remaining group-theoretic properties are standard.
The standard $S$-arithmetic lattice theorem realizes $G$ diagonally as a
lattice in
\[
 \prod_{v\mid\infty}\SL_n(K_v)
 \times\prod_{\mathfrak q\in\Sigma}\SL_n(K_{\mathfrak q})
\]
\cite{BorelSerre}.  The local-field theorem for $\SL_n$ and
lattice inheritance give property $(T)$
\cite[Theorems~1.4.15 and~1.7.1]{BHV}.  The finiteness theorem of
Borel--Serre gives finite presentability \cite{BorelSerre}, and linearity is
immediate.

To see residual finiteness directly, let $g\neq1$ and choose a nonzero entry
$x\in\OS$ of $g-I_n$.  A nonzero element of the Dedekind domain $\OS$ lies
in only finitely many maximal ideals, so one can choose a nonzero prime ideal
$\mathfrak q\notin\Sigma$ with $x\notin\mathfrak q\OS$.  Reduction modulo
$\mathfrak q$ maps $G$ to the finite group $\SL_n(\OK/\mathfrak q)$ and
does not kill $g-I_n$.
\end{proof}

In the rational case the class-group step disappears.  If $p\in S$, then
$\alpha=p$ and the displayed calculation reads
\[
 \operatorname{diag}(pI_{n-1},p^{-(n-1)})
 \begin{pmatrix}A&x\\0&1\end{pmatrix}
 \operatorname{diag}(p^{-1}I_{n-1},p^{n-1})
 =\begin{pmatrix}A&p^n x\\0&1\end{pmatrix}.
\]
The affine property-$(T)$ group and its finite-index self-compressions also
appear in \cite[Proposition~1.1]{deCornulier}.

\section{Finite-dimensional approximation}\label{sec:completion}

For a group $G$ and $x\in\C[G]$, set
\[
 \|x\|_{\mathrm{fd}}
 =\sup\{\|\pi(x)\|:\pi\text{ is a finite-dimensional unitary
 representation of }G\}.
\]
Finite-dimensional exotic group completions have been studied previously;
see \cite{Wiersma,RuanWiersma}.  The compression defect makes both quotient
kernels below explicit.

\begin{proposition}\label{prop:completion}
Let $G$ be residually finite and satisfy Theorem~\ref{thm:criterion}.  Then
\[
 \|x\|_r\leq\|x\|_{\mathrm{fd}}\leq\|x\|_{\max}
 \qquad(x\in\C[G]),
\]
and $\|\cdot\|_{\mathrm{fd}}$ is a $C^*$-norm on $\C[G]$.  Its completion
$\Cfd G$ is residually finite-dimensional and stably finite.  Both canonical
quotients
\[
 \Cmax G\longrightarrow\Cfd G\longrightarrow\Cr G
\]
are proper: the first kernel contains
\[
 d=p_\Gamma-u_{t^{-1}}p_\Gamma u_t,
\]
and the second kernel contains the image of $p_\Gamma$.
\end{proposition}

\begin{proof}
Let $0\neq x=\sum_{g\in F}a_gu_g\in\C[G]$ with $F$ finite.  Residual
finiteness gives a finite quotient of $G$ that is injective on $F$.  The
image of $x$ in the group algebra of this quotient is nonzero, and the
regular representation of the quotient is faithful.  Thus
$\|x\|_{\mathrm{fd}}>0$.

To compare with the reduced norm, let $\xi\in\ell^2(G)$ be finitely
supported.  There is a finite quotient $\theta:G\to Q$ that is injective on
\[
 F\cup\operatorname{supp}(\xi)\cup F\operatorname{supp}(\xi).
\]
Transporting $\xi$ to $\ell^2(Q)$ then preserves both its norm and the norm
of its image under $x$.  Taking the supremum over such $\xi$ gives
$\|x\|_r\leq\|x\|_{\mathrm{fd}}$.

The product of all finite-dimensional representations realizes $\Cfd G$ as
a $C^*$-subalgebra of a product of matrix algebras, so it is residually
finite-dimensional and stably finite.  Proposition~\ref{prop:defect} gives
$0\neq d\in\ker(\Cmax G\to\Cfd G)$.  The image of $p_\Gamma$ in
$\Cfd G$ is nonzero because the augmentation character sends it to $1$,
whereas its image in $\Cr G$ is zero.  Hence the second quotient is proper.
\end{proof}

Thus every group in Theorem~\ref{thm:sarithmetic} has a canonical
residually finite-dimensional, stably finite exotic completion strictly
between its maximal and reduced group algebras.

The construction of Proposition~\ref{prop:kazhdan} also gives finite-support
witnesses for the norm loss.  Namely,
\[
 a_j=L^j-u_{t^{-1}}L^ju_t\in\C[G]
\]
satisfies
\[
 \|a_j-d\|_{\max}\leq2r^j,
 \qquad
 \|a_j\|_{\mathrm{fd}}\leq2r^j,
 \qquad
 \|\lambda_{G/\Gamma}(a_j)\|\geq1-2r^j.
\]

\begin{thebibliography}{99}

\bibitem{Bekka}
M.~B. Bekka,
\emph{On the full $C^*$-algebras of arithmetic groups and the congruence
subgroup problem},
Forum Math. \textbf{11} (1999), no.~6, 705--715.

\bibitem{BekkaSuper}
M.~B. Bekka,
\emph{Operator-algebraic superrigidity for $\mathrm{SL}_n(\mathbb Z)$,
$n\geq3$},
Invent. Math. \textbf{169} (2007), no.~2, 401--425.

\bibitem{BHV}
B.~Bekka, P.~de la Harpe, and A.~Valette,
\emph{Kazhdan's Property $(T)$},
New Mathematical Monographs, vol.~11, Cambridge University Press,
Cambridge, 2008.

\bibitem{BlackadarKirchberg}
B.~Blackadar and E.~Kirchberg,
\emph{Generalized inductive limits of finite-dimensional $C^*$-algebras},
Math. Ann. \textbf{307} (1997), no.~3, 343--380.

\bibitem{BorelHarishChandra}
A.~Borel and Harish-Chandra,
\emph{Arithmetic subgroups of algebraic groups},
Ann. of Math. (2) \textbf{75} (1962), 485--535.

\bibitem{BorelSerre}
A.~Borel and J.-P.~Serre,
\emph{Cohomologie d'immeubles et de groupes $S$-arithm\'etiques},
Topology \textbf{15} (1976), no.~3, 211--232.

\bibitem{CarrionDadarlatEckhardt}
J.~R. Carri\'on, M.~Dadarlat, and C.~Eckhardt,
\emph{On groups with quasidiagonal $C^*$-algebras},
J. Funct. Anal. \textbf{265} (2013), no.~1, 135--152.

\bibitem{ChoiQDF}
Y.~Choi,
\emph{A note on some group $C^*$-algebras which are quasi-directly finite},
arXiv:1003.1650v2 (2010).

\bibitem{Coburn}
L.~A. Coburn,
\emph{The $C^*$-algebra generated by an isometry},
Bull. Amer. Math. Soc. \textbf{73} (1967), 722--726.

\bibitem{deCornulier}
Y.~de Cornulier,
\emph{Finitely presentable, non-Hopfian groups with Kazhdan's property
$(T)$ and infinite outer automorphism group},
Proc. Amer. Math. Soc. \textbf{135} (2007), no.~4, 951--959.

\bibitem{DrutuNowak}
C.~Dru\c tu and P.~W. Nowak,
\emph{Kazhdan projections, random walks and ergodic theorems},
J. Reine Angew. Math. \textbf{754} (2019), 49--86.

\bibitem{LacaLarsen}
M.~Laca and N.~S. Larsen,
\emph{Hecke algebras of semidirect products},
Proc. Amer. Math. Soc. \textbf{131} (2003), no.~7, 2189--2199;
erratum, \textbf{133} (2005), no.~4, 1255--1256.

\bibitem{MilhojRordam}
H.~O. Milh\o j and M.~R\o rdam,
\emph{Around traces and quasitraces},
accepted for M\"unster J. Math.; arXiv:2309.17412v2.

\bibitem{RuanWiersma}
Z.-J. Ruan and M.~Wiersma,
\emph{On exotic group $C^*$-algebras},
J. Funct. Anal. \textbf{271} (2016), no.~2, 437--453.

\bibitem{Shulman}
T.~Shulman,
\emph{The MF property for amalgamated free products},
arXiv:2603.13564v2 (2026).

\bibitem{Wassermann}
S.~Wassermann,
\emph{$C^*$-algebras associated with groups with Kazhdan's property $T$},
Ann. of Math. (2) \textbf{134} (1991), no.~2, 423--431.

\bibitem{Wiersma}
M.~Wiersma,
\emph{Constructions of exotic group $C^*$-algebras},
Illinois J. Math. \textbf{60} (2016), no.~3, 655--667.

\end{thebibliography}

\end{document}
